Write \(q=p+1\).  Fix smooth functions \(\phi,\eta:[0,\infty)\to[0,1]\) such that \(\phi=1\) near zero, \(\phi=0\) on \([1,\infty)\), all transition endpoints are flat, \(\eta(u)=0\) for \(u\leq1\), and \(\eta(u)=1\) for \(u\geq2\).  The radial function may be chosen as a smooth function of squared radius, so its pullback to \(\mathbb R^2\) is smooth at the origin.  Choose constants in the order: a small slope \(a>0\), a large cutoff constant \(C\), and a large smoothness constant \(A\).  Put
\[
 w=A\Delta^{1/q}.
\]
Place \(x_j=(s_j,0)\) on \([-1/2,1/2]\times\{0\}\) with separation at least \(4w\).  Elementary interval packing gives
\[
 M\asymp w^{-1}\geq c\Delta^{-1/q},
\]
and the disks \(S_j=B_2(x_j,w)\) are disjoint, contained in \(\mathcal X=[-3,3]^2\), and avoid the corners of \(\mathcal B=\partial([-1,1]\times[0,2])\).

For \(r>0\), define
\[
 b_{\Delta,w}(r)={2\Delta\over ar}\phi(r/w)\eta(r/(C\Delta)),
 \qquad b_{\Delta,w}(0)=0.
\]
For \(x\in S_j\), write \(x-x_j=r_j(\cos\alpha_j,\sin\alpha_j)\), and interpret every summand below as zero outside its disk.  For \(\omega\in\{0,1\}^M\), let
\[
 f_\omega(x)={1\over36}\left[1+\sum_{j=1}^M\omega_jb_{\Delta,w}(r_j)\cos\alpha_j\right],
\]
\[
 \mu_{0,\omega}(x)={1\over2},
 \qquad
 \mu_{1,\omega}(x)={1\over2}+ax_1-\Delta\sum_{j=1}^M\omega_j\phi(r_j/w).
\]
Conditional on \(X=x\), take \(Y(0)\) and \(Y(1)\) conditionally independent Bernoulli variables with these respective success probabilities.  Treatment and observation obey the fixed deterministic assignment and consistency equations in the class definition.

We verify the law-class clauses one by one.  The density perturbation vanishes on \(r\leq C\Delta\) and smoothly before \(r=w\), so the apparent factor \(1/r\) creates no singularity and \(f_\omega\) is continuous on the square.  Whenever it is nonzero,
\[
 |b_{\Delta,w}(r)|\leq {2\over aC}.
\]
Choose \(C\) so large that this is at most \(1/4\).  Because the disks are disjoint,
\[
 {1\over48}\leq f_\omega(x)\leq {5\over144}.
\]
Moreover, polar integration on every full disk gives
\[
 \int_0^w b_{\Delta,w}(r)r\,dr\int_0^{2\pi}\cos\alpha\,d\alpha=0.
\]
Since the square has area \(36\), \(\int_{\mathcal X}f_\omega=1\).  Thus the density is normalized, positive on the entire fixed support, continuous, and satisfies the required envelope once \(L_0\geq48\).

For every derivative order \(0\leq k\leq q\), scaling and disjointness give
\[
 \left\|D^k\left[\Delta\sum_j\omega_j\phi(\lVert\,\cdot-x_j\rVert_2/w)\right]\right\|_\infty
 \leq C_k\Delta w^{-k}=C_kA^{-k}\Delta^{1-k/q}.
\]
At \(k=q\) the right side is \(C_qA^{-q}\), and for \(k<q\) it tends to zero.  Choosing \(A\) large verifies both the derivative suprema through order \(p\) and the Lipschitz bounds for every such derivative.  The formula extends smoothly past the support because all disks are strictly interior and the affine term is global.  Taking \(a\) and then \(\Delta\) small puts both Bernoulli probabilities in \([1/4,3/4]\).  Hence their conditional variances lie in \([3/16,1/4]\), their raw conditional \((2+\nu)\)-moments are at most one for every \(\nu\geq2\), and all the variance, moment, and regression clauses hold.  Independent replication gives the required sampling condition for every \(n\).

The assignment rectangle has perimeter eight and is at Euclidean distance at least one from the support boundary.  Its boundary is therefore a compact positive-length rectifiable curve in the support interior.  At a smooth edge or a corner and for all sufficiently small \(h\), each arm contains a circular sector centered at the boundary point with opening angle at least \(\pi/2\).  Since the kernel equals one on the unit ball, scaling such a sector by \(h^{-1}\) gives
\[
 {1\over h^2}\int_{\mathcal A_t}K_\square(\{(2t-1)\lVert z-x\rVert_2\}/h)\,dz\geq {\pi\over8}.
\]
The corresponding radius-\(s\) arc has length between \((\pi/2)s\) and \(2\pi s\); multiplication by the density envelope proves that the Hausdorff denominator is positive and finite at edges and corners alike.

For the Gram condition, use polar coordinates in one of those sectors.  For a coefficient vector \(v\in\mathbb R^{p+1}\) and the sign \(\epsilon_t=2t-1\), the density lower bound gives
\[
 v^\top\Psi_{t,P_\omega,x}(h)v
 \geq c\int_0^1\{v^\top r_p(\epsilon_tu)\}^2u\,du.
\]
The integral is strictly positive for every nonzero \(v\), because a nonzero polynomial cannot vanish on an interval.  It is continuous on the unit sphere, which is compact, so its minimum there is positive.  This yields a lower eigenvalue depending only on \(p\) and the fixed density and sector constants.  Enlarging \(L_0(p)\) verifies the displayed uniform Gram bound.  Finally, the Euclidean-ball VC proof in \(\mathrm{lem:euclidean\text{-}balls\text{-}vc}\) applies to this fixed metric, including after restriction to the square.  Thus every A1/A2 and Theorem~2 clause holds for every vertex.

We next establish locality and target separation.  Each disk perturbation integrates to zero, so
\[
 P_\omega(X\in S_j)={|S_j|\over36}={\pi w^2\over36}=\rho
\]
for every \(j\) and \(\omega\).  Disjoint supports show that the joint law restricted to \(S_j\) depends only on \(\omega_j\), while all perturbations vanish outside their union.  At \(x_j\),
\[
 \tau_{P_\omega}(x_j)=\mu_{1,\omega}(x_j)-\mu_{0,\omega}(x_j)=as_j-\Delta\omega_j,
\]
so adjacent target values differ by exactly \(\Delta\).

It remains to calculate the compressed information.  Inside \(S_j\), positive signed radius is the upper semicircle \(0\leq\alpha\leq\pi\) and negative signed radius is the lower semicircle.  The identities
\[
 \int_0^\pi\cos\alpha\,d\alpha=0,
 \qquad
 \int_0^\pi\cos^2\alpha\,d\alpha={\pi\over2}
\]
show that both signed-radius marginals are bit-independent.  Conditional on the disk, their density is \(r/w^2\) on each of \((0,w)\) and \((-w,0)\).  On the control half, the observed Bernoulli mean is always \(1/2\).  On the treatment half, angular averaging gives, for \(r>0\),
\[
 q_{0,j}(r)={1\over2}+as_j,
\]
and
\[
 q_{1,j}(r)-q_{0,j}(r)
 =-\Delta\phi(r/w)+{ar\over2}b_{\Delta,w}(r)
 =-\Delta\phi(r/w)\{1-\eta(r/(C\Delta))\}.
\]
Thus the complete conditional outcome--signed-distance laws agree unless \(0<r<2C\Delta\).  Their Bernoulli probabilities remain in \([1/4,3/4]\), so in either KL orientation
\[
 \operatorname{kl}(u,v)\leq { (u-v)^2\over v(1-v)}\leq C_1\Delta^2.
\]
Disintegration over the common signed-radius marginal yields
\[
 \operatorname{KL}(Q_{j,0},Q_{j,1})
 \leq C_1\Delta^2\int_0^{2C\Delta}{r\over w^2}\,dr
 \leq C_0{\Delta^4\over w^2},
\]
and reversing the adjacent vertices gives the same bound.  This proves every assertion.